\subsection{Whole local norm polynomials and nilpotent lengths}
\label{subsec:local-norm-polynomials}

These are whole determinant polynomial identities, obtained from actual
nilpotent elements of the polynomial tensor algebra. Equality of scalar
polynomial values over a finite original field is not used.

\begin{lemma}[Determinant norm of a nilpotent scalar perturbation]
\label{spectrum:nilpotent-scalar-determinant}
\lean{EnvelopeProperty.algebraNorm_eq_pow_of_isNilpotent_sub}
\leanok
Let $R$ be an integral domain and $E$ a finite free commutative $R$-algebra.
If $b-\operatorname{map}(c)$ is nilpotent for $b\in E$ and $c\in R$, then
\[
 N_{E/R}(b)=c^{\operatorname{rank}_R E}.
\]
\end{lemma}
\begin{proof}
Multiplication by $b-\operatorname{map}(c)$ is nilpotent. Its
characteristic polynomial is therefore $X^d$, where $d$ is the actual
finite free rank. Multiplication by $b$ differs from this nilpotent
operator by the scalar operator $c$. The characteristic polynomial
shift formula identifies its polynomial with $(X-c)^d$. Its signed
constant coefficient, which is its determinant, is $c^d$.
This is an equality of actual algebra norms over the domain, not an
estimate on values or a reducedness assumption.
\end{proof}

\begin{lemma}[Actual residual scalar of a local algebra]
\label{spectrum:algclosed-local-residue-scalar}
\lean{EnvelopeProperty.localResidueScalar,
EnvelopeProperty.localResidueScalar_residue,
EnvelopeProperty.isNilpotent_sub_localResidueScalar,
EnvelopeProperty.finrank_eq_localLength}
\leanok
\uses{spectrum:local-algebra-norm-length}
Let $K$ be algebraically closed and $E$ a finite local commutative
$K$-algebra. The actual coefficient map $K\to\kappa(E)$ is an isomorphism;
its inverse after the residue map defines $\rho:E\to K$. Every
$b-\operatorname{map}(\rho(b))$ is nilpotent, and
$\dim_K E=\ell_E(E)$.
\end{lemma}
\begin{proof}
The actual residue field is finite, hence integral, over $K$.
Algebraic closedness makes its coefficient map bijective. Compose the
residue map with the genuine inverse algebra isomorphism to define
$\rho$; the residue of its coefficient inclusion recovers the residue
of $b$ on every element.

An element with residue zero belongs to the maximal ideal. In an
Artinian local ring that ideal is the Jacobson radical and also the
nilradical, so the displayed difference is genuinely nilpotent.
Finally, the residue isomorphism gives residue degree one. The
already proved identity between dimension, local length and residue
degree therefore identifies the dimension with the actual local length.
\end{proof}

\begin{theorem}[Whole local norm form with nilpotent multiplicity]
\label{spectrum:algclosed-local-norm-form}
\lean{EnvelopeProperty.normForm_eq_localResidueScalar_pow,
EnvelopeProperty.wgaussValuation_normForm_of_local_algClosed}
\leanok
\uses{spectrum:nilpotent-scalar-determinant,
spectrum:algclosed-local-residue-scalar,spectrum:whole-norm-form-base-change}
For $E,K,\rho$ as above and a finite family $b_j\in E$, the actual
norm polynomial satisfies
\[
 N_{E/K}\Bigl(\sum_j X_jb_j\Bigr)
 =\Bigl(\sum_j\rho(b_j)X_j\Bigr)^{\ell_E(E)}.
\]
For any real ultrametric valuation $v$ on $K$ and any radii $r_j\ge0$,
\[
 \left\|N_{E/K}\Bigl(\sum_j X_jb_j\Bigr)\right\|_r
 =\left(\max_jv(\rho(b_j))r_j\right)^{\ell_E(E)}.
\]
\end{theorem}
\begin{proof}
Put $P=K[X_j]$. In the actual tensor algebra $P\otimes_K E$, subtract
the coefficient inclusion of the residual scalar polynomial from
$z=\sum_jX_j\otimes b_j$. Tensor balancing and the actual coefficient
map identify the difference with
\[
 \sum_j (X_j\otimes1)
       \bigl(1\otimes(b_j-\operatorname{map}(\rho(b_j)))\bigr).
\]
Every parenthesized element is nilpotent by the preceding lemma and
its image under the tensor inclusion remains nilpotent. The nilradical
is an ideal, so this entire finite sum is genuinely nilpotent.

The polynomial ring $P$ is an integral domain and the tensor algebra
is finite free over it. Apply the exact scalar-perturbation determinant
formula to $z$ and the residual scalar polynomial. Tensor base change
preserves the finite free rank, which equals $\ell_E(E)$ by the preceding
lemma. By definition this determinant is the original norm polynomial,
proving equality of the entire polynomials. Weighted Gauss
multiplicativity and the exact linear-form evaluation yield the
second identity. This proof retains nilpotent layers through the
local length and invokes no polynomial interpolation, separability,
reducedness, or complete norm on $K$.
\end{proof}
